\begin{pastexam}[label={px:ust-115-ee-ds-07}]{臺聯大電機類 115 學年度 第 7 題}
（單選；2 分）

How do we use a stack structure when evaluating postfix expressions (e.g., 3 4 + 5 * 7 / 2-)?
\begin{enumerate}[label=(\Alph*)]
\item When encountering an operand, push; when encountering an operator, pop two numbers. Perform
the operation on the two numbers and then push the result.
\item When encountering an operator, use Push; when encountering an operand, use Pop.
\item Push all operands and operators encountered, and then calculate again at the end.
\item This requires using columns, not stacks.
\end{enumerate}
\end{pastexam}

\begin{pastsoln}
答案是 (A)。

後序式求值（資結書第~6.3.2~節、虛擬碼~6.7）：由左往右讀，遇到運算元就 push；遇到運算子就 pop
兩次，先 pop 出來的是\textbf{右}運算元、後 pop 的是左運算元，算完把結果 push 回去。 讀完時堆疊裡
剩下的一個數就是答案。 追蹤題目的例子：
\begin{center}
\begin{tabular}{lll}
\toprule
讀到 & 動作 & 堆疊（左邊是底） \\
\midrule
3 & push & 3 \\
4 & push & 3, 4 \\
$+$ & pop 4、3，算 $3 + 4$ & 7 \\
5 & push & 7, 5 \\
$*$ & pop 5、7，算 $7 \times 5$ & 35 \\
7 & push & 35, 7 \\
$/$ & pop 7、35，算 $35 / 7$ & 5 \\
2 & push & 5, 2 \\
$-$ & pop 2、5，算 $5 - 2$ & 3 \\
\bottomrule
\end{tabular}
\end{center}
結果是 $3$，也就是中序式 $(3 + 4) \times 5 / 7 - 2$。

逐項檢查：
\begin{itemize}
\item (A) 正是上面的做法。 對。
\item (B) 反過來了：運算子不必進堆疊，運算元才要進堆疊等著被用。 錯。
\item (C) 全部 push 完再算，等於沒有利用後序式「讀到運算子時兩個運算元剛好在頂端」的性質；而且
算的時候還是要用上面的規則。 錯。
\item (D) 後序求值的標準工具就是堆疊。 「columns」大概是「佇列」之類的誤譯，不論哪一種，都不是
必須的。 錯。
\end{itemize}
\end{pastsoln}
